\documentclass[11pt]{article}
\usepackage[margin=1in]{geometry}
\usepackage[T1]{fontenc}
\usepackage{xcolor}
\usepackage[utf8]{inputenc}
\usepackage{booktabs,tabularx,enumitem,xurl,hyperref,microtype,array}
\hypersetup{colorlinks=true,linkcolor=blue!50!black,urlcolor=blue!50!black}
\newcolumntype{L}{>{\raggedright\arraybackslash}X}
\setlength{\parindent}{0pt}\setlength{\parskip}{6pt}
\DeclareUnicodeCharacter{2011}{-}
\title{Positioning and Pain in Credit Markets: a Literature Review}
\author{Prepared for the PositioningWork project}
\date{1 November 2026}
\begin{document}
\maketitle

\section*{The answer}
\textbf{The literature supports flow and forced-selling effects in credit, but nobody has shown a positioning or Pain signal that predicts credit out of sample.} We found no such test for corporate bonds on RFQ, credit ETFs, credit futures or CDX. So our failed tests neither confirm nor contradict a published result. Five points matter most.

\begin{enumerate}
\item \textbf{A loss alone does not force selling.} Effects appear when a loss meets a binding limit, such as margin, leverage, regulatory capital or fund redemptions (Siriwardane 2019; Kang et al.\ 2020; Fang 2024).
\item \textbf{Who holds the position decides the sign.} Levered accounts sell as limits approach (Bian et al.\ 2018). Insurers sell winners and keep losers (Ellul et al.\ 2015). One Pain number across all holders can cancel itself out.
\item \textbf{Our zero of 53 cannot separate decay from low power.} Equity signals keep about half their performance after publication (McLean and Pontiff 2016). With about 14 independent observations, a signal that kept half its power would pass one test only about 17\% of the time. That arithmetic comes from the review, not from a paper.
\item \textbf{Our cost-basis method has no published validation.} No paper we found estimates cost basis from open interest or CFTC data. The standard method also retires old lots with trading volume (Grinblatt and Han 2005), and ours does not.
\item \textbf{Credit data has its own traps.} Price errors and one-sided filters inflate most of 108 published bond factors (Dickerson et al.). One widely used factor paper, Bai, Bali and Wen (2019), has been retracted.
\end{enumerate}

\section{Corporate bonds and RFQ}
\textbf{Some client flow predicts returns, and some reverses.} Czech and Pinter (2020) use UK regulatory data with both parties named. Flow from well-connected clients predicts returns out to 30 days, with a top-minus-bottom spread of 17.1 bp. Flow from poorly connected clients reverses by 30 days.

\textbf{Measured losses change trading.} Fang (2024) compares insurers trading the same bond at the same time, using book values. Insurers with larger unrealized losses provide less liquidity against fund fire sales. The price impact of flow-induced selling is $-0.406$ percentage points of yield spread per 1\% of amount outstanding. One standard deviation more peer loss adds $-0.420$. These numbers come from a working paper.

\textbf{Selling pressure usually reverses.} This holds for insurer forced selling, sell herding and flow-driven fund pressure. Effects are largest in illiquid and high-yield bonds. That is also where price noise and costs are highest.

\textbf{The open question is RFQ itself.} We found no paper that tests RFQ imbalance as a return or spread predictor. Published RFQ work studies costs and search. No public dataset we found gives client type at trade level.

\textbf{What breaks flow measures in bonds:} basket trades, quasi-dealers on all-to-all venues, clients who are really supplying liquidity, and repeated inquiries for one order. The TRACE error rate is also material: Dick-Nielsen (2009) reports 7.7\% of reports are errors.

\section{Credit ETFs}
\textbf{Share count records dealer work more than investor views.} For US corporate bond ETFs, the top three authorised participants do 92\% of creations and redemptions, and bank participants do 90\% (Cai et al.\ 2023). Baskets are small and they change (Shim and Todorov 2021; Koont et al.\ 2025). So our null result on share count is plausible. No paper tests it directly.

\textbf{Short interest mixes plumbing with bets.} Evans et al.\ (2026) show that market makers sell ETF shares short and create them later. That operational shorting does not predict one-month returns. Their return test covers equity ETFs only. Xiao (2022) finds higher short interest after redemption days in corporate bond ETFs.

\textbf{Whether flows push bond prices is contested.}
\begin{itemize}
\item Dannhauser and Hoseinzade (2021): yes in the 2013 Taper Tantrum, reversing after about four months.
\item Rossi et al.\ (2026 draft): no significant reaction.
\item Agapova et al.\ (2025): ETFs act as a buffer, not as demand pressure.
\end{itemize}
The studies differ in period and flow measure. The effect may exist only when dealers are constrained. That reconciliation is the review's reading, and no paper tests it directly.

\textbf{What the literature points to instead.} In order of evidence: split short interest into operational and directional parts, use bond-level basket data, condition on dealer capacity, and test price-based signals such as premiums. All are the review's reading. None is an established recipe for credit.

\section{CDX and credit futures}
\textbf{Public data cannot show who holds CDX.} Rule 17 CFR 43.4 bars disclosing swap parties in public trade data. Papers on holder types use confidential regulatory data.

\textbf{The Pain idea has its strongest support here.} Siriwardane (2019) finds that margin payments by protection sellers explain 12\% of the variation in weekly CDS spread changes. It works because it measures losses at the actual holders. Public data cannot rebuild it.

\textbf{But levered CDX holders were not visibly the forced side in March 2020.} A CFTC staff paper shows hedge funds buying about \$80 billion of protection net into 20 March (Coughlan et al.\ 2022). Bond investors may have been the forced sellers. So do not assume that CDX holders are always the ones forced out. This also tempers our earlier view that CDX is the natural home for Pain.

\textbf{Rolls distort positions.} Positions spike around the March and September rolls. Average one-month rises reached 22\% over the six rolls before March 2020 (Coughlan et al.\ 2022). Any flow signal must adjust for them.

\textbf{Credit futures have almost no history.} The CFTC reports Bloomberg HY, IG and IG Duration-Hedged futures, with 24, 26 and 8 reports. Only 20 to 25 traders report, and the top four IG traders held 76.8\% of longs on 27 October 2026. High yield has not appeared since 15 September 2026. Cboe IBHY and IBIG do not appear in the CFTC reports. We found no academic study of credit futures positioning.

\section{Positioning, Pain and test design}
\textbf{The closest published analogue to Pain} is Kang et al.\ (2020), on 26 commodity futures. After a large hedger capital loss, the return impact of a given position change more than doubles, from 8.9 to 21.2 bp. Those numbers come from a 2016 draft. Their loss measure needs no entry prices, so it is a simple benchmark for our ledger.

\textbf{Three cost-basis methods exist:} weighting past prices by turnover (Grinblatt and Han 2005), building from fund holdings (Frazzini 2006), and reading book values from insurer filings (Fang 2024). Splitting gains from losses also matters, because the loss side can behave differently (An 2016).

\textbf{Higher-power tests are possible.} Use weekly non-overlapping returns, since long horizons add little independent evidence (Boudoukh et al.\ 2008). Pool across contracts with robust inference. Test interactions, such as whether impact grows when losses and limits coincide.

\section{Where the evidence conflicts}
\small
\begin{tabularx}{\linewidth}{@{}p{0.27\linewidth}LL@{}}
\toprule
Question & Positions & Likely reason \\
\midrule
Do ETF flows move bond prices? & Yes and reverses (Dannhauser and Hoseinzade). No (Rossi et al.). Buffer (Agapova et al.). & Stress window against broader samples \\
Do fund outflows cause fire sales? & Little pressure (Choi et al.\ 2020). Real effects (Falato et al.; Fang). & Same-issuer controls and sample \\
Do positions predict returns? & Yes in risk-premium tests (De Roon et al.; Kang et al.). Almost nothing in Granger tests (Sanders et al.). & The level mixes opposing demands \\
Is bond momentum real? & Yes (Jostova et al.). Mostly noise and filters (Dickerson et al.). & The later work tests the earlier method \\
Who was forced in March 2020? & Bond investors (Haddad et al.). Hedge funds were fast net buyers of CDX protection (Coughlan et al.). & These can fit together \\
\bottomrule
\end{tabularx}
\normalsize

\section{What this means for your four markets}
These are our inferences, not findings.
\begin{enumerate}
\item \textbf{TRACE and RFQ.} Look for short-term pressure that reverses, and measure it against dealer capacity. Use spread-based entry levels. Your anonymous, single-dealer RFQ data cannot show client breadth, which is the split that mattered in Czech and Pinter.
\item \textbf{Credit ETFs.} Drop share count as a position. Test basket data, directional short interest and premiums.
\item \textbf{CDX.} Use proxies such as index skew and cleared totals, and accept that they do not show holders. Adjust for rolls.
\item \textbf{Credit futures.} Monitor only. A few firms can move the series.
\item \textbf{Measured holders.} Insurer book values in NAIC filings are the one real cost basis in credit. Access terms are unchecked.
\end{enumerate}

\section{Gaps the literature leaves open}
\begin{enumerate}
\item Any out-of-sample test of positions forecasting CDX spreads.
\item Any test of RFQ imbalance against bond returns or spreads.
\item Any test of bond ETF share count or short interest as positioning.
\item Any cost-basis estimate from open interest, checked against true positions.
\item Any post-publication test of CFTC positioning signals.
\item Any corporate bond test of the disposition effect outside insurers.
\end{enumerate}

\section{How this review was done, and its limits}
Five parallel AI-assisted searches covered credit ETFs, corporate bonds and RFQ, CDX and futures, positioning theory, and local documents. The search strings and exclusions are in the theme files in this folder. The review also covers 42 local documents, summarised in \texttt{E\_local\_documents.md}. Local documents gave no test on credit, and \texttt{ContraFlow.pdf} reports no evidence.

\textbf{Checks.} All 196 distinct DOIs in the theme files resolve in Crossref. I checked first authors automatically on 72 ranked entries that have a standard heading. I inspected the two flags, and both were false positives from parsing. The check did not cover group B's ranked entries, because its format differs. A DOI that resolves shows only that the paper exists. It does not confirm a stated finding.

\textbf{Limits.}
\begin{itemize}
\item Many entries were read at abstract level or in a working-paper draft. Each entry below states what was read, so check the source before relying on a number.
\item Some scholarly search services were rate-limited or failed, so recent papers may be missing. SSRN blocked direct access for some ETF entries.
\item DTCC, CME, Cboe and BIS data pages could not be fully confirmed.
\item This is not a PRISMA systematic review. No meta-analysis was run.
\item AI tools found, read and summarised these sources. A person should check any citation before external use.
\end{itemize}

\appendix
\section{Annotated bibliography}
\small
\subsection*{A. Credit ETFs}
\begin{enumerate}[leftmargin=*,itemsep=3pt]
\item Dannhauser, C. D., and Hoseinzade, S. (2021). The Unintended Consequences of Corporate Bond ETFs: Evidence from the Taper Tantrum. Review of Financial Studies. \nolinkurl{10.1093/rfs/hhab031}. \emph{Tier 1. Read: abstract.} \textbf{For us:} The flow effect is real in one stress episode and it reverses. Test outflows in stress, with a long enough horizon, and do not expect an all-weather signal.
\item Evans, R. B., Moussawi, R., Pagano, M. S., and Sedunov, J. (2026). Operational Shorting and ETF Liquidity Provision. Journal of Financial Economics. Earlier title: ETF Short Interest and Failures-to-Deliver: Naked Short-Selling or Operational Shorting? (SSRN 2017). \nolinkurl{10.1016/j.jfineco.2026.104241}. \emph{Tier 1. Read: abstract, introduction and the return test sections of the working paper (Wharton Jacobs Levy copy).} \textbf{For us:} This is the most direct warning for us. Raw ETF short interest mixes dealer plumbing with real bets. Separate them before testing. The bond ETF case is untested here.
\item Brown, D. C., Davies, S. W., and Ringgenberg, M. C. (2021). ETF Arbitrage, Non-Fundamental Demand, and Return Predictability. Review of Finance, 25(4), 937-972. \nolinkurl{10.1093/rof/rfaa027}. \emph{Tier 1. Read: abstract.} \textbf{For us:} This is the main positive case for flows. The sign is contrarian and the measure is the change in shares, not the level. Check that our tests allow for that. The result may not carry over to credit.
\item Koont, N., Ma, Y., Pastor, L., and Zeng, Y. (2025). Steering a Ship in Illiquid Waters: Active Management of Passive Funds. Review of Financial Studies. Also NBER WP 30039. \nolinkurl{10.1093/rfs/hhaf034}. \emph{Tier 1. Read: abstract, plus the data and liquidity sections of the NBER PDF.} \textbf{For us:} A share change does not map to a pro-rata trade in bonds. Bond-level basket data is closer to true pressure than fund-level share count.
\item Shim, J. J., and Todorov, K. (2021). ETFs, Illiquid Assets, and Fire Sales. Working paper. A search listing also shows it as BIS Working Paper 975, which I could not open. \nolinkurl{10.2139/ssrn.3886881}. \emph{Tier 2. Read: abstract, introduction and basket sections of the 9 May 2022 draft (RSM copy).} \textbf{For us:} Dealer inventory matters more than share count for ETF discounts. A discount can mean dealers are protecting their book, not that investors are bearish.
\item Todorov, K. (2021). The anatomy of bond ETF arbitrage. BIS Quarterly Review, March 2021. \nolinkurl{https://www.bis.org/publ/qtrpdf/r\_qt2103d.htm}. \emph{Tier 3. Read: full text.} \textbf{For us:} Arbitrage in credit ETFs works through a narrow channel. A fund-level flow can touch few bonds. Note: this is the Todorov piece I could confirm. I found no separate Todorov paper titled on ETF discounts.
\item Pan, K., and Zeng, Y. (2017). ETF Arbitrage under Liquidity Mismatch. ESRB Working Paper 59. Also SSRN. \nolinkurl{10.2139/ssrn.2895478}. \emph{Tier 2. Read: abstract and data sections.} \textbf{For us:} Creations and redemptions partly reflect the AP's own bond book. That breaks the reading of share changes as end-investor positioning.
\item Liu, D. D., and Ge, L. (2026). From ETF Baskets to Bond Prices: Creation, Redemption, and Price Pressure. SSRN. \nolinkurl{10.2139/ssrn.7551378}. \emph{Tier 2. Read: search excerpt only.} \textbf{For us:} This points to a bond-level basket imbalance as the better flow measure. It is very new and not peer reviewed. I could not check whether it is predictive or same-period.
\item Agapova, A., Kaprielyan, M., and Volkov, N. (2025). ETFs and the price volatility of underlying bonds. Financial Review, 60(3), 667-700. \nolinkurl{10.1111/fire.12426}. \emph{Tier 1. Read: abstract.} \textbf{For us:} This is a partial no-effect result that fits ours. High-yield flows did nothing for volatility, and the daily sign runs against simple flow pressure.
\item Rossi, M., Ruan, Y., and Xia, S. (2026 draft, first posted 2023). Real Effects of Corporate Bond ETF Ownership. Working paper. \nolinkurl{10.2139/ssrn.4425722}. \emph{Tier 2. Read: abstract and introduction of the 2 April 2026 draft.} \textbf{For us:} This is a direct no-effect result on ETF outflows and bond returns. It conflicts with entry 1. Ownership moves spreads in level, but flows may not.
\item Xiao, H. (2022). The Economics of ETF Redemptions. Working paper. SSRN title adds: Liquidity and Fire Sale Risks. \nolinkurl{10.2139/ssrn.4096222}. \emph{Tier 2. Read: abstract, introduction and the regression sections of the 29 December 2022 draft.} \textbf{For us:} This is the closest match to our sign pattern. Falling shares go with rising short interest in credit ETFs. The author reads the short interest as arbitrage activity, not a view. One-day timing depends on how share data is dated, so check our own lag.
\item Hao, X., Kim, D., Sul, H. K., and Wang, Q. (2019). Underlying Bond Return Predictability by ETF Returns. SSRN. \nolinkurl{10.2139/ssrn.3570309}. \emph{Tier 2. Read: abstract.} \textbf{For us:} The ETF price carries the signal, not the quantity. Predictability is stronger where flows are smaller, which argues against flows as the carrier.
\item Haddad, V., Moreira, A., and Muir, T. (2021). When Selling Becomes Viral: Disruptions in Debt Markets in the COVID-19 Crisis and the Fed's Response. Review of Financial Studies. Also NBER WP 27168. \nolinkurl{10.1093/rfs/hhaa145}. \emph{Tier 1. Read: abstract, plus a keyword scan of the ETF passages in the NBER PDF.} \textbf{For us:} In March 2020 selling was about raising cash, not credit views. Discounts were largest in the safest funds, so a raw discount is not a clean credit signal. Policy events can dominate any position measure.
\item Huang, A. G., O'Hara, M., Xue, J., and Zhou, X. (2025). Beyond Creation and Redemption: Principal Trading by Bond ETFs. SSRN. \nolinkurl{10.2139/ssrn.5760891}. \emph{Tier 2. Read: search excerpt only.} \textbf{For us:} Funds trade bonds a lot with no change in share count. Share count therefore misses much of the bond-market footprint.
\item Dannhauser, C. D., and Karmaziene, E. (2023). The Dealer Warehouse: Corporate Bond ETFs. SSRN. \nolinkurl{10.2139/ssrn.4613818}. \emph{Tier 2. Read: search excerpt only.} \textbf{For us:} A rise in shares can mean dealers are parking bonds they want to sell. That is the opposite of investors adding credit risk.
\item Reilly, C. (2022). The Hidden Cost of Corporate Bond ETFs. SSRN. \nolinkurl{10.2139/ssrn.4147882}. \emph{Tier 2. Read: abstract as quoted on Quantpedia.} \textbf{For us:} The informed signal sits in which bonds APs deliver, not in how many shares they create.
\item Raddatz K., C. E. (2025). Authorized participants' regulatory constraints and limits to ETF arbitrage during market turmoil: Evidence from the dash-for-cash episode. Journal of Banking and Finance, 179. \nolinkurl{10.1016/j.jbankfin.2025.107499}. \emph{Tier 1. Read: abstract.} \textbf{For us:} The meaning of a share change depends on AP balance sheets. Condition on AP capital in stress.
\item Dinger, J., Muller, M., Rzeznik, A., and Uhrig-Homburg, M. (2025). Same Same But Different: The Risk Profile of Corporate Bond ETFs. SSRN. \nolinkurl{10.2139/ssrn.5193084}. \emph{Tier 2. Read: abstract and introduction of the December 2025 draft.} \textbf{For us:} ETF prices track dealer capacity more than investor demand. An AP or dealer capital measure may do more than a position measure.
\item Chen, J. (2023). ETF Shorting Activity, Return Predictability and Information Efficiency of Underlying Securities. SSRN. \nolinkurl{10.2139/ssrn.4527115}. \emph{Tier 2. Read: thesis abstract and the sample passage.} \textbf{For us:} This is the bullish case for ETF short interest, but only in equities. It conflicts with entry 2 until operational shorting is removed. Do not assume it carries over to credit.
\item Ben-Rephael, A., Choi, J., and Goldstein, I. (2021). Mutual fund flows and fluctuations in credit and business cycles. Journal of Financial Economics, 139(1), 84-108. \nolinkurl{10.1016/j.jfineco.2020.07.004}. \emph{Tier 1. Read: abstract.} \textbf{For us:} This is the strongest flow-based credit signal I found. It uses exchanges within fund families. My reading is that this filters out some liquidity trading and dealer activity, which the abstract does not say. It is a slow-moving credit cycle signal, not a trading signal.
\item Fulkerson, J. A., Jordan, S. D., and Riley, T. B. (2013). Predictability in Bond ETF Returns. Journal of Fixed Income, 23(3). \nolinkurl{10.3905/jfi.2013.23.3.050}. \emph{Tier 1. Read: abstract.} \textbf{For us:} The premium or discount is a more promising ETF-level signal than shares. Stale NAVs and trading costs may explain part of this. That caution is mine, not the authors'.
\item Dannhauser, C. D. (2017). The impact of innovation: Evidence from corporate bond exchange-traded funds (ETFs). Journal of Financial Economics, 125(3), 537-560. \nolinkurl{10.1016/j.jfineco.2017.06.002}. \emph{Tier 1. Read: abstract.} \textbf{For us:} ETF ownership shifts bond valuations over the long run. That is a slow level effect, not a signal for next-period returns. A no-liquidity-gain result too.
\item Finnerty, J. D., Reisel, N., and Zhong, X. (2025). ETFs, Creation and Redemption Processes, and Bond Liquidity. Journal of Financial and Quantitative Analysis, 60(4), 1891-1924. \nolinkurl{10.1017/S0022109024000346}. \emph{Tier 1. Read: abstract.} \textbf{For us:} Redemptions are not simply bearish pressure. They can improve bond liquidity. That weakens a one-directional reading of share falls.
\item Marta, T. (2021 draft, first posted 2019). Fixed Income ETFs, Bond Liquidity, and Stressed Markets. Working paper. SSRN title: Fixed Income ETFs and Bond Liquidity. \nolinkurl{10.2139/ssrn.3350519}. \emph{Tier 2. Read: abstract and sample passages of the August 2021 draft.} \textbf{For us:} This is a no-harm result. It conflicts with entries 1, 22 and with Meli and Todorova. ETF volume may tell us more than share count.
\item Cai, F., Chuan, G., Henry, K., Shin, C., and Tuzun, T. (2023). New Insights from N-CEN: Liquidity Management at Open-End Funds and Primary Market Concentration of ETFs. FEDS Notes, 11 January 2023. \nolinkurl{https://www.federalreserve.gov/econres/notes/feds-notes/new-insights-from-n-cen-liquidity-management-at-open-end-funds-and-pmc-of-etfs-20230111.html}. \emph{Tier 3. Read: full text.} \textbf{For us:} A handful of dealer banks produce the share count series. It tells us about those few balance sheets. It is not a wide sample of investors.
\end{enumerate}
\subsection*{B. Corporate bonds and RFQ}
\begin{enumerate}[leftmargin=*,itemsep=3pt]
\item Czech, R. and Pinter, G. (2020). Informed Trading and the Dynamics of Client-Dealer Connections in Corporate Bond Markets. Bank of England Staff Working Paper 895. \nolinkurl{10.2139/ssrn.3744325}. \emph{Tier 2. Read: part text of the December 2020 version.} \textbf{For us:} split RFQ flow by client breadth. Do not treat all client flow as one signal. Count how many dealers a client uses, which a single dealer sees only in part.
\item Fang, C. (2024, draft April 2025). Unrealized Trading Gains. Working paper. \nolinkurl{10.2139/ssrn.4973722}. \emph{Tier 2. Read: part text of the April 2025 draft.} \textbf{For us:} this is the best template for "pain". It measures loss against an entry level and links it to trading and price impact. For insurers the entry level is book value, not last trade price. Losses on similar bonds can matter more than losses on the bond itself.
\item Geilen, M. (2026). Corporate Bond Dealers and Price Discovery. Working paper. \nolinkurl{10.2139/ssrn.6781998}. \emph{Tier 2. Read: abstract.} \textbf{For us:} dealer interest and client interest are different signals. Test RFQ flow alongside axes where we can get axes. This paper is new and has not been through peer review.
\item Dickerson, A., Robotti, C. and Rossetti, G. (2026). The Corporate Bond Factor Replication Crisis. Working paper. \nolinkurl{10.2139/ssrn.6088966}. \emph{Tier 2. Read: abstract of the 2026 paper.} \textbf{For us:} never build the entry level and the forward return from the same noisy print. Skip a period between signal and return. Use benchmark prices for marks. Backtest with the corrected data as well as our own.
\item Ellul, A., Jotikasthira, C., Lundblad, C. and Wang, Y. (2015). Is Historical Cost Accounting a Panacea? Market Stress, Incentive Distortions, and Gains Trading. Journal of Finance 70(6), 2489-2538. \nolinkurl{10.1111/jofi.12357}. \emph{Tier 1. Read: abstract.} \textbf{For us:} stress can lead an investor to sell winners, not losers. So a pain index built only on losers can miss the selling. Track gains as well as losses by client type.
\item Jurkatis, S., Schrimpf, A., Todorov, K. and Vause, N. (2023). Relationship Discounts in Corporate Bond Trading. Bank of England Staff Working Paper 1,049. \nolinkurl{10.2139/ssrn.4301416}. \emph{Tier 2. Read: part text.} \textbf{For us:} the informed-client score is directly usable on RFQ data. Client level also includes a price effect, so control for relationship when we use the level as an entry price.
\item Kargar, M., Lester, B., Plante, S. and Weill, P.-O. (2023). Sequential Search for Corporate Bonds. NBER Working Paper 31904. \nolinkurl{10.3386/w31904}. \emph{Tier 2. Read: part text.} \textbf{For us:} this is the closest published use of RFQ records. A single order can appear as several inquiries, so we must merge repeats before we count flow. Failed inquiries carry unfilled demand that TRACE cannot see.
\item Cai, F., Han, S., Li, D. and Li, Y. (2019). Institutional Herding and Its Price Impact: Evidence from the Corporate Bond Market. Journal of Financial Economics 131(1), 139-167. \nolinkurl{10.1016/j.jfineco.2018.07.012}. \emph{Tier 1. Read: abstract and the data section of the FEDS 2016-091 version.} \textbf{For us:} expect asymmetry. Build buy and sell crowding as two signals. Crowded selling in illiquid high yield looks like a contrarian setup, not a momentum one.
\item Anand, A., Jotikasthira, C. and Venkataraman, K. (2021). Mutual Fund Trading Style and Bond Market Fragility. Review of Financial Studies 34(6), 2993-3044. \nolinkurl{10.1093/rfs/hhaa120}. \emph{Tier 1. Read: abstract.} \textbf{For us:} classify clients by behaviour, not just sector. A client's history of supplying or demanding liquidity tells us how to read its flow.
\item Friewald, N. and Nagler, F. (2024). Dealer Inventory and the Cross-Section of Corporate Bond Returns. Economics Letters 239, 111710. \nolinkurl{10.1016/j.econlet.2024.111710}. \emph{Tier 1. Read: abstract.} \textbf{For us:} customer selling becomes dealer inventory, so one-sided flow should earn a premium to whoever absorbs it. This is the main control for any flow signal. Weekly returns of this size may not survive costs, and entry 4 applies.
\item Jiang, H. and Sun, Z. (2013 draft). Understanding the Illiquidity of Corporate Bonds: The Arrival of Public News. Working paper. \nolinkurl{https://cba.lmu.edu/media/lmucollegeofbusinessadministration/undergrad/05\_Sun.pdf}. \emph{Tier 2. Read: part text.} \textbf{For us:} condition on size and on public news. Flow on news days carries less private information. Treat this as weak evidence, because it is an old draft with no journal record that I could find.
\item Goyal, A., Nozawa, Y. and Qiu, Y. (FINRA lists October 2025, SSRN record 2026). On the Drivers of Corporate Bond Lending. Working paper. \nolinkurl{10.2139/ssrn.7392021}. \emph{Tier 2. Read: both abstract versions.} \textbf{For us:} do not import equity short-interest logic. High bond lending can mean customers are buying, which reverses the usual reading.
\item Hendershott, T., Kozhan, R. and Raman, V. (Crossref date 2018). Short Selling and Price Discovery in Corporate Bonds. Journal of Financial and Quantitative Analysis 55(1), 77-115. \nolinkurl{10.1017/s0022109018001539}. \emph{Tier 1. Read: abstract.} \textbf{For us:} read with entry 12. Aggregate flow signals to the issuer as well as the bond. Include past stock returns as a control.
\item Jacobsen, S., Venkataraman, K. and Xu, D. (2025, revised January 2026). Illiquidity Meets Intelligence: AI-Driven Price Discovery in Corporate Bonds. Working paper. \nolinkurl{10.2139/ssrn.5297667}. \emph{Tier 2. Read: abstract on the FINRA page.} \textbf{For us:} RFQ information already leaks into a vendor reference price. Any RFQ signal must beat CP+ or a similar mark. CP+ may also be a better entry and mark price than the last trade.
\item Comerton-Forde, C., Ford, B., Foucault, T. and Jurkatis, S. (2025). Investors as a Liquidity Backstop in Corporate Bond Markets. HEC Paris Research Paper. \nolinkurl{10.2139/ssrn.5229934}. \emph{Tier 2. Read: abstract on the FINRA page.} \textbf{For us:} a client's sell can be liquidity it supplies to a dealer, not its own view. A signed flow signal should try to separate the two.
\item Ellul, A., Jotikasthira, C. and Lundblad, C. (2011). Regulatory Pressure and Fire Sales in the Corporate Bond Market. Journal of Financial Economics 101(3), 596-620. \nolinkurl{10.1016/j.jfineco.2011.03.020}. \emph{Tier 1. Read: abstract, published and SSRN versions.} \textbf{For us:} the event trigger is rating plus holder constraint, not the downgrade alone. Forced selling implies a reversal that a positioning signal could trade.
\item Choi, J., Hoseinzade, S., Shin, S. and Tehranian, H. (2020). Corporate Bond Mutual Funds and Asset Fire Sales. Journal of Financial Economics 138(2), 432-457. \nolinkurl{10.1016/j.jfineco.2020.05.006}. \emph{Tier 1. Read: abstract.} \textbf{For us:} fund outflow alone is a weak bond-level signal. Selling also lands on the liquid bonds, not the weakest. Compare within issuer so that we do not mistake credit news for pressure.
\item Falato, A., Hortacsu, A., Li, D. and Shin, C. (2021). Fire-Sale Spillovers in Debt Markets. Journal of Finance 76(6), 3055-3102. \nolinkurl{10.1111/jofi.13078}. \emph{Tier 1. Read: abstract of the SSRN version, 10.2139/ssrn.2887822.} \textbf{For us:} measure overlap across holders. Risk comes from who else holds the bond, and not only from the client we see.
\item Goldstein, I., Jiang, H. and Ng, D. (2017). Investor Flows and Fragility in Corporate Bond Funds. Journal of Financial Economics 126(3), 592-613. \nolinkurl{10.1016/j.jfineco.2016.11.007}. \emph{Tier 1. Read: abstract, published and SSRN versions.} \textbf{For us:} under-water funds are the ones whose end investors leave. Pain should feed through fund performance and flow, and it should be non-linear on the downside.
\item Coppola, A. (2025). In Safe Hands: The Financial and Real Impact of Investor Composition over the Credit Cycle. Review of Financial Studies 38(8), 2275-2325. \nolinkurl{10.1093/rfs/hhaf017}. \emph{Tier 1. Read: abstract and data section of the draft.} \textbf{For us:} scale any pain or crowding measure by holder type. The same selling means more in a fund-held bond than in an insurer-held one.
\item Chaudhary, M., Fu, J. Z. and Li, J. (2026). Corporate Bond Multipliers: Substitutes Matter. Review of Financial Studies, forthcoming. \nolinkurl{10.1093/rfs/hhag082}. \emph{Tier 1. Read: SSRN abstract only.} \textbf{For us:} test flow at sector, rating and maturity bucket level, not only by CUSIP. Single-bond pressure may just move to close substitutes.
\item O'Hara, M. and Zhou, X. (2021). The Electronic Evolution of Corporate Bond Dealers. Journal of Financial Economics 140(2), 368-390. \nolinkurl{10.1016/j.jfineco.2021.01.001}. \emph{Tier 1. Read: abstract.} \textbf{For us:} RFQ data is a selected slice of flow. Measure platform share by bond and period. Otherwise growth in platform use looks like growth in demand.
\item Hendershott, T., Livdan, D. and Schurhoff, N. (2021). All-to-All Liquidity in Corporate Bonds. Swiss Finance Institute Research Paper 21-43. \nolinkurl{https://papers.ssrn.com/sol3/papers.cfm?abstract\_id=3895270}. \emph{Tier 2. Read: part text.} \textbf{For us:} "client" on a platform may be a quasi-dealer. Filter these liquidity providers out of positioning measures, or their inventory trades will look like views.
\item Li, J., O'Hara, M., Rapp, A. and Zhou, X. (2023, draft February 2025). Bond Market Illiquidity: Is Portfolio Trading the Solution? Working paper. \nolinkurl{10.2139/ssrn.4495516}. \emph{Tier 2. Read: part text.} \textbf{For us:} basket trades do not express a view on each bond. Tag them and exclude them, or give them less weight. Their price is also not a clean entry level per bond.
\item Meli, J. and Todorova, Z. (2022). Portfolio Trading in Corporate Bond Markets. Working paper. \nolinkurl{10.2139/ssrn.4295453}. \emph{Tier 2. Read: part text of the draft and the published abstract.} \textbf{For us:} dealers hedge baskets with ETFs, so client flow may show up in ETF flow. Link this to theme work on credit ETFs.
\item Bessembinder, H., Jacobsen, S., Maxwell, W. and Venkataraman, K. (2018). Capital Commitment and Illiquidity in Corporate Bonds. Journal of Finance 73(4), 1615-1661. \nolinkurl{10.1111/jofi.12694}. \emph{Tier 1. Read: abstract.} \textbf{For us:} signal strength may vary by regime. Dealers now absorb less, so one-sided client flow may need to find other clients. Do not assume a constant relation across 2006 to today.
\item Dick-Nielsen, J. and Rossi, M. (2019). The Cost of Immediacy for Corporate Bonds. Review of Financial Studies 32(1), 1-41. \nolinkurl{10.1093/rfs/hhy080}. \emph{Tier 1. Read: abstract.} \textbf{For us:} index exclusions give a clean test of uninformed flow. Use them to separate the price effect of pressure from the effect of information in our indicator.
\item O'Hara, M. and Zhou, X. (2021). Anatomy of a Liquidity Crisis: Corporate Bonds in the COVID-19 Crisis. Journal of Financial Economics 142(1), 46-68. \nolinkurl{10.1016/j.jfineco.2021.05.052}. \emph{Tier 1. Read: abstract.} \textbf{For us:} normal rules break in stress. Larger trades can stop being cheaper, and RFQ flow can shift towards liquid bonds. The indicator needs a stress regime, and its level prices will be poor marks there.
\item Feldhutter, P. (2012). The Same Bond at Different Prices: Identifying Search Frictions and Selling Pressures. Review of Financial Studies 25(4), 1155-1206. \nolinkurl{10.1093/rfs/hhr093}. \emph{Tier 1. Read: abstract.} \textbf{For us:} this gives a pain measure that RFQ size and level data can compute directly. Entry 28 warns that size pricing can invert in a crisis, which is what the measure looks for.
\item Jostova, G., Nikolova, S., Philipov, A. and Stahel, C. (2013). Momentum in Corporate Bond Returns. Review of Financial Studies 26(7), 1649-1693. \nolinkurl{10.1093/rfs/hht022}. \emph{Tier 1. Read: SSRN abstract, 10.2139/ssrn.1651853.} \textbf{For us:} use as the classic control. Entry 4 and the 2023 Dickerson paper dispute that such return-based profits survive noise correction, so treat the size with care.
\end{enumerate}
\subsection*{C. CDX and credit futures}
\begin{enumerate}[leftmargin=*,itemsep=3pt]
\item John Coughlan, Madison Lau, Alexei Orlov. "Single-Name and Index CDS Dynamics during the Market Stress of 2020." CFTC Office of the Chief Economist, Staff Papers and Reports 2022-007, October 2022. \nolinkurl{https://www.cftc.gov/sites/default/files/2022-10/CDS\%20white\%20paper\%20October\%202022\_ada.pdf}. \emph{Tier 3. Read: abstract, sections 1 to 4.2, and hedge fund passages.} \textbf{For us:} Gross notional is a poor stress gauge because dealer intermediation dominates it. Any position change near 20 March or 20 September must be adjusted for the roll. A text search of the paper found no discussion of forced unwinds, so this episode does not by itself show forced exits.
\item Emil N. Siriwardane. "Limited Investment Capital and Credit Spreads." Journal of Finance 74(5), 2303-2347. \nolinkurl{10.1111/jofi.12777}. \emph{Tier 1. Read: abstract only.} \textbf{For us:} This is the closest published precedent for a CDS pain indicator. It works because it measures losses at the actual holders. Public data cannot rebuild it, so we would have to proxy holder losses. I did not read the paper's concentration statistics.
\item Agostino Capponi, Wan-Schwin Allen Cheng, Stefano Giglio, Richard Haynes. "The Collateral Rule: Evidence from the Credit Default Swap Market." Journal of Monetary Economics 126, 58-86. \nolinkurl{10.1016/j.jmoneco.2021.12.003}. \emph{Tier 1. Read: abstract of the CFTC preliminary draft (https://www.cftc.gov/idc/groups/public/@economicanalysis/documents/file/oce\_collateral\_rule.pdf).} \textbf{For us:} A forced-exit signal should be nonlinear. Small spread moves should score near zero, and large moves should score sharply higher. This is the strongest available support for our levered-holder hypothesis, though the paper studies margins and not exits.
\item Xinjie Wang, Zhaodong Zhong. "Dealer Inventory, Pricing, and Liquidity in the OTC Derivatives Markets: Evidence from Index CDSs." Journal of Financial Markets 57. \nolinkurl{10.1016/j.finmar.2020.100617}. \emph{Tier 1. Read: abstract only.} \textbf{For us:} This is the most direct index evidence on positions and prices. The abstract describes association, not forecasting. Do not cite it as proof that positions predict spreads.
\item Nina Boyarchenko, Pooja Gupta, Nick Steele, Jacqueline Yen. "Trends in Credit Basis Spreads." Federal Reserve Bank of New York Economic Policy Review 24(2). \nolinkurl{https://www.newyorkfed.org/research/epr/2018/epr\_2018\_trends-in-credit-basis-spreads\_boyarchenko}. \emph{Tier 2. Read: full relevant sections.} \textbf{For us:} Index skew measures arbitrage capacity, not just mispricing. Real-money CDX selling is also part of the positioning story, so hedge funds are not the only side to watch.
\item Benjamin Junge, Anders B. Trolle. "Liquidity Risk in Credit Default Swap Markets." Swiss Finance Institute Research Paper 13-65. \nolinkurl{10.2139/ssrn.2371374}. \emph{Tier 2. Read: abstract and introduction.} \textbf{For us:} This gives a defined skew construction that we can adapt. It needs index and constituent CDS quotes, which are usually licensed. Its high persistence suggests using changes or surprises.
\item Pierre Collin-Dufresne, Benjamin Junge, Anders B. Trolle. "Market Structure and Transaction Costs of Index CDSs." Journal of Finance. Crossref year 2020. \nolinkurl{10.1111/jofi.12953}. \emph{Tier 1. Read: published abstract, plus parts of the September 2017 draft (https://www.aeaweb.org/conference/2018/preliminary/paper/HA3dNy8B).} \textbf{For us:} Client flow appears to carry information, but public tapes do not label who traded. Size caps also censor the largest trades, which are the ones most likely to matter for forced exits.
\item Nina Boyarchenko, Anna M. Costello, Or Shachar. "The Long and Short of It: The Post-Crisis Corporate CDS Market." Federal Reserve Bank of New York Staff Report 879, revised August 2019. \nolinkurl{https://www.newyorkfed.org/medialibrary/media/research/staff\_reports/sr879.pdf}. \emph{Tier 2. Read: abstract and selected sections.} \textbf{For us:} Dealer sign can flip over regimes, so a fixed reading of "dealer long" is unsafe. Positions across single-name, index, options, and tranches need reading together. U.S.-only data miss much of the other side.
\item Jaewon Choi, Or Shachar, Sean Seunghun Shin. "Dealer Liquidity Provision and the Breakdown of the Law of One Price: Evidence from the CDS-Bond Basis." Management Science. \nolinkurl{10.1287/mnsc.2017.3011}. \emph{Tier 1. Read: abstract.} \textbf{For us:} This is the cleanest evidence in the theme that arbitrageur exit moves prices. It concerns bonds in a crisis, so transfer to CDX is a hypothesis.
\item Andrea L. Eisfeldt, Bernard Herskovic, Sriram Rajan, Emil N. Siriwardane. "OTC Intermediaries." Review of Financial Studies. \nolinkurl{10.1093/rfs/hhac062}. \emph{Tier 1. Read: abstract.} \textbf{For us:} Dealer capacity is a state variable. A positioning signal may need conditioning on dealer stress, although this is a model calibration and not a test of a signal.
\item Mark Mitchell, Todd Pulvino. "Arbitrage Crashes and the Speed of Capital." Journal of Financial Economics 104(3), 469-490. \nolinkurl{10.1016/j.jfineco.2011.09.002}. \emph{Tier 1. Read: first part of the abstract.} \textbf{For us:} Funding withdrawal, not just spread loss, can force exits. A pain indicator should allow for a funding input such as hedge fund credit availability.
\item Nina Boyarchenko, Thomas M. Eisenbach, Pooja Gupta, Or Shachar, Peter Van Tassel. "Bank-Intermediated Arbitrage." Federal Reserve Bank of New York Staff Report 858, revised July 2020. \nolinkurl{10.2139/ssrn.3200041}. \emph{Tier 2. Read: abstract and part of the trade definitions.} \textbf{For us:} Skew can stay wide without being a clean opportunity. High skew may reflect low capacity, and that needs separating from a positioning signal.
\item Martin Oehmke, Ádám Zawadowski. "The Anatomy of the CDS Market." Review of Financial Studies. \nolinkurl{10.1093/rfs/hhw068}. \emph{Tier 1. Read: published abstract, plus parts of the 2013 draft (https://www.aeaweb.org/conference/2013/retrieve.php?pdfid=59).} \textbf{For us:} Net notional is usable as a size measure but does not reveal direction by holder. These are single-name results, so index use is an extension.
\item Jennie Bai, Pierre Collin-Dufresne. "The CDS-Bond Basis." Financial Management. \nolinkurl{10.1111/fima.12252}. \emph{Tier 1. Read: abstract.} \textbf{For us:} Friction proxies can condition basis signals. These are bond-level results, not index results.
\item Iñaki Aldasoro, Torsten Ehlers. "The Credit Default Swap Market: What a Difference a Decade Makes." BIS Quarterly Review, June 2018. \nolinkurl{https://www.bis.org/publications/qr-201806/credit-default-swap-market-what-difference-decade-makes}. \emph{Tier 3. Read: key takeaways and hedge fund passage.} \textbf{For us:} BIS is a rare public hedge fund counterparty view. It is semiannual, global, and dealer-reported, so it suits regime context and not a trading signal.
\item Lynn Riggs, Esen Onur, David Reiffen, Haoxiang Zhu. "Swap Trading after Dodd-Frank: Evidence from Index CDS." Journal of Financial Economics. \nolinkurl{10.1016/j.jfineco.2020.03.008}. \emph{Tier 1. Read: abstract.} \textbf{For us:} Large holders may hide intent and split execution with few dealers. Public volume may therefore understate stress until after the move.
\item Yee Cheng Loon, Zhaodong Zhong. "Does Dodd-Frank Affect OTC Transaction Costs and Liquidity? Evidence from Real-Time CDS Trade Reports." Journal of Financial Economics 119(3). \nolinkurl{10.1016/j.jfineco.2016.01.019}. \emph{Tier 1. Read: abstract.} \textbf{For us:} This confirms that public trade tapes support academic liquidity measures. That gives us a stress and liquidity input even without holder identity.
\item Wenxin Du, Salil Gadgil, Michael B. Gordy, Clara Vega. "Counterparty Risk and Counterparty Choice in the Credit Default Swap Market." Management Science. \nolinkurl{10.1287/mnsc.2023.4870}. \emph{Tier 1. Read: abstract.} \textbf{For us:} Stress may first show up in clearing and counterparty shifts rather than in price. Public cleared-share series may be a low-cost stress input, which is untested.
\item Mark Paddrik, Stathis Tompaidis. "Market-Making Costs and Liquidity: Evidence from CDS Markets." Office of Financial Research Working Paper 19-01. \nolinkurl{10.2139/ssrn.3351288}. \emph{Tier 2. Read: abstract.} \textbf{For us:} Aggregate dealer positions, the only kind public data could ever show, may miss the pricing channel. This limits what a TFF dealer series can tell us.
\item Mark Paddrik, Stathis Tompaidis. "Intermediation Networks and Derivative Market Liquidity: Evidence from Credit Default Swap Markets." Review of Finance. Earlier version: OFR Working Paper 24-01. \nolinkurl{10.1093/rof/rfag023}. \emph{Tier 1. Read: abstract.} \textbf{For us:} Network structure matters, but public data do not reveal it. Treat this as a known omitted variable.
\item Richard Haynes, Lihong McPhail. "The Liquidity of Credit Default Index Swap Networks." CFTC Office of the Chief Economist, June 2017. \nolinkurl{10.2139/ssrn.3378629}. \emph{Tier 3. Read: abstract.} \textbf{For us:} Execution cost differs by holder, so forced exits by narrowly connected holders may cost more. That is an untested extension.
\item Emil N. Siriwardane, Adi Sunderam, Jonathan Wallen. "Segmented Arbitrage." Journal of Finance. \nolinkurl{10.1111/jofi.13469}. \emph{Tier 1. Read: abstract.} \textbf{For us:} Do not import a general intermediary stress index and assume it explains CDX. CDX capacity may be local to its own participants.
\item Valentin Haddad, Alan Moreira, Tyler Muir. "When Selling Becomes Viral: Disruptions in Debt Markets in the COVID-19 Crisis and the Fed's Response." Review of Financial Studies. \nolinkurl{10.1093/rfs/hhaa145}. \emph{Tier 1. Read: abstract.} \textbf{For us:} Stress can be led by cash-bond sellers rather than CDX holders. A CDX signal must allow for selling pressure that starts elsewhere.
\item Andrea L. Eisfeldt, Bernard Herskovic, Shuo Liu. "Interdealer Price Dispersion and Intermediary Capacity." NBER Working Paper 32998. \nolinkurl{10.3386/w32998}. \emph{Tier 2. Read: abstract.} \textbf{For us:} Dispersion is a candidate capacity signal. We have not confirmed that it can be built from public CDX data.
\item Cboe Futures Exchange, LLC. Rule Certification, Submission Number CFE-2026-011, letter to the CFTC, 16 March 2026. \nolinkurl{https://cdn.cboe.com/resources/regulation/rule\_filings/pending/2026/26-011-Credit-Futures-Position-Limits.pdf}. \emph{Tier 3. Read: product, statistics, and rationale sections.} \textbf{For us:} Cboe has real open interest but no public holder split that I could find. Position limits are also a structural cap on forced-exit size, and they changed in 2026.
\end{enumerate}
\subsection*{D. Positioning, pain and test design}
\begin{enumerate}[leftmargin=*,itemsep=3pt]
\item Grinblatt, M., and B. Han. "Prospect theory, mental accounting, and momentum." Journal of Financial Economics 78(2), 311-339. \nolinkurl{10.1016/j.jfineco.2004.10.006}. \emph{Tier 1. Read: RePEc abstract of the journal version.} \textbf{For us:} this is the standard estimator. Its removal rule matches our pro-rata rule, but it retires lots using volume, not only net changes.
\item An, L. "Asset Pricing When Traders Sell Extreme Winners and Losers." Review of Financial Studies 29(3), 823-861. \nolinkurl{10.1093/rfs/hhv060}. \emph{Tier 1. Read: abstract.} \textbf{For us:} pain may not be one-directional. Test large gains and large losses as separate signals, not a single signed score.
\item Frazzini, A. "The Disposition Effect and Underreaction to News." Journal of Finance 61(4), 2017-2046. \nolinkurl{10.1111/j.1540-6261.2006.00896.x}. \emph{Tier 1. Read: abstract and text (author PDF).} \textbf{For us:} holdings beat volume as a source where they exist. The removal rule may matter little, which we can test because our ledger already has three rules.
\item Riley, C., B. Summers, and D. Duxbury. "Capital Gains Overhang with a Dynamic Reference Point." Management Science 66(10), 4726-4745. \nolinkurl{10.1287/mnsc.2019.3404}. \emph{Tier 1. Read: abstract.} \textbf{For us:} entry price may not be the only anchor. Recent highs and lows in spread could be added and tested.
\item Ben-David, I., and D. Hirshleifer. "Are Investors Really Reluctant to Realize Their Losses? Trading Responses to Past Returns and the Disposition Effect." Review of Financial Studies 25(8), 2485-2532. \nolinkurl{10.1093/rfs/hhs077}. \emph{Tier 1. Read: abstract.} \textbf{For us:} do not assume a sharp trigger at the entry price. "Under water" may not be a meaningful threshold.
\item Cici, G. "The Prevalence of the Disposition Effect in Mutual Funds' Trades." Journal of Financial and Quantitative Analysis 47(4), 795-820. \nolinkurl{10.1017/s0022109012000348}. \emph{Tier 1. Read: abstract.} \textbf{For us:} professional investors may not behave as the retail evidence implies. This effect itself appears to have weakened over time.
\item Hincapié-Salazar, J., and D. Agudelo. "Is the disposition effect in bonds as strong as in stocks? Evidence from an emerging market." Global Finance Journal 46, 100508. \nolinkurl{10.1016/j.gfj.2019.100508}. \emph{Tier 1. Read: abstract.} \textbf{For us:} this is the closest direct test I found, and it points against us. The holders who dominate US credit may not show this bias.
\item Fang, C. "Unrealized Trading Gains." SSRN working paper. \nolinkurl{10.2139/ssrn.4973722}. \emph{Tier 2. Read: abstract and text of the April 2025 version.} \textbf{For us:} this is the best template we have. It uses a measured cost basis, not an estimate. Losses here reduce buying, which is a different channel from forced selling.
\item Ellul, A., C. Jotikasthira, C. Lundblad, and Y. Wang. "Is Historical Cost Accounting a Panacea? Market Stress, Incentive Distortions, and Gains Trading." Journal of Finance 70(6), 2489-2538. \nolinkurl{10.1111/jofi.12357}. \emph{Tier 1. Read: journal abstract and the abstract of the December 2011 draft.} \textbf{For us:} under stress, some holders sell their winners and keep their losers. A pain indicator that assumes losers get sold could have the wrong sign for insurers.
\item Kang, W., K. G. Rouwenhorst, and K. Tang. "A Tale of Two Premiums: The Role of Hedgers and Speculators in Commodity Futures Markets." Journal of Finance 75(1), 377-417. \nolinkurl{10.1111/jofi.12845}. \emph{Tier 1. Read: journal abstract.} \textbf{For us:} this is the closest published analogue to a pain signal. Its weekly loss measure is a simple benchmark for our ledger. The level of positioning mixes two opposing signals.
\item Bian, J., Z. He, K. Shue, and H. Zhou. "Leverage-Induced Fire Sales and Stock Market Crashes." NBER Working Paper 25040. \nolinkurl{10.3386/w25040}. \emph{Tier 2. Read: abstract.} \textbf{For us:} the trigger is distance to a limit, not the size of the loss. The expected return signature is a decline, then a reversal.
\item Brunnermeier, M., and L. Pedersen. "Market Liquidity and Funding Liquidity." Review of Financial Studies 22(6), 2201-2238. \nolinkurl{10.1093/rfs/hhn098}. \emph{Tier 1. Read: abstract.} \textbf{For us:} pain should matter most when funding and volatility are stressed. It implies a conditional test, not an unconditional one.
\item Siriwardane, E. "Limited Investment Capital and Credit Spreads." Journal of Finance 74(5), 2303-2347. \nolinkurl{10.1111/jofi.12777}. \emph{Tier 1. Read: abstract.} \textbf{For us:} this is the strongest direct evidence for credit. Mark-to-market losses on positions do move spreads. Note that it uses real positions and a weekly horizon, and its result is about spread changes.
\item Goldstein, I., H. Jiang, and D. Ng. "Investor flows and fragility in corporate bond funds." Journal of Financial Economics 126(3), 592-613. \nolinkurl{10.1016/j.jfineco.2016.11.007}. \emph{Tier 1. Read: RePEc abstract.} \textbf{For us:} fund losses transmit to selling through investors, not through a manager's entry price. Fund performance and illiquidity may be better inputs than cost basis.
\item Choi, J., S. Hoseinzade, S. Shin, and H. Tehranian. "Corporate bond mutual funds and asset fire sales." Journal of Financial Economics 138(2), 432-457. \nolinkurl{10.1016/j.jfineco.2020.05.006}. \emph{Tier 1. Read: RePEc abstract.} \textbf{For us:} this conflicts with the forced-selling story. Issuer news may account for apparent fire sales, so a pain test needs to control for it.
\item De Roon, F., T. Nijman, and C. Veld. "Hedging Pressure Effects in Futures Markets." Journal of Finance 55(3), 1437-1456. \nolinkurl{10.1111/0022-1082.00253}. \emph{Tier 1. Read: abstract.} \textbf{For us:} pressure from related markets may matter. An index-level signal could use rates or equity positions as well.
\item Sanders, D., S. Irwin, and R. Merrin. "Smart Money? The Forecasting Ability of CFTC Large Traders." Proceedings of the NCCC-134 Conference, Chicago. \nolinkurl{https://farmdoc.illinois.edu/assets/meetings/nccc134/conf\_2007/pdf/confp01-07.pdf}. \emph{Tier 2. Read: abstract and introduction.} \textbf{For us:} a negative result on plain positioning, predating our test by over a decade. Prices may lead our positions, so check the lead-lag direction first.
\item Hong, H., and M. Yogo. "What does futures market interest tell us about the macroeconomy and asset prices?" Journal of Financial Economics 105(3), 473-490. \nolinkurl{10.1016/j.jfineco.2012.04.005}. \emph{Tier 1. Read: RePEc abstract.} \textbf{For us:} open interest may carry information without any entry price, so it makes a useful baseline. The evidence is weakest outside commodities and has no out-of-sample support.
\item Huang, D., J. Li, L. Wang, and G. Zhou. "Time series momentum: Is it there?" Journal of Financial Economics 135(3), 774-794. \nolinkurl{10.1016/j.jfineco.2019.08.004}. \emph{Tier 1. Read: abstract.} \textbf{For us:} pooling raises power, but pooled t-statistics can mislead. Any pooled pain test should use a bootstrap.
\item Lou, D., and C. Polk. "Comomentum: Inferring Arbitrage Activity from Return Correlations." Review of Financial Studies 35(7), 3272-3302. \nolinkurl{10.1093/rfs/hhab117}. \emph{Tier 1. Read: abstract.} \textbf{For us:} crowding can be inferred from returns alone, with no position data. Such a measure could condition our signal.
\item Barroso, P., R. Edelen, and P. Karehnke. "Crowding and Tail Risk in Momentum Returns." Journal of Financial and Quantitative Analysis 57(4), 1313-1342. \nolinkurl{10.1017/s0022109021000624}. \emph{Tier 1. Read: abstract.} \textbf{For us:} high positioning alone may not predict a crash. This supports testing crowding only together with a trigger.
\item McLean, R. D., and J. Pontiff. "Does Academic Research Destroy Stock Return Predictability?" Journal of Finance 71(1), 5-32. \nolinkurl{10.1111/jofi.12365}. \emph{Tier 1. Read: abstract.} \textbf{For us:} plan on keeping about half of a backtest's return. A sell-side backtest may deserve a bigger haircut.
\item Falck, A., A. Rej, and D. Thesmar. "When do systematic strategies decay?" Quantitative Finance 22(11), 1955-1969. \nolinkurl{10.1080/14697688.2022.2098810}. \emph{Tier 1. Read: short journal abstract.} \textbf{For us:} our signal has many construction steps, so these results point to above-average decay risk. Simpler versions deserve a test.
\item Hou, K., C. Xue, and L. Zhang. "Replicating Anomalies." Review of Financial Studies 33(5), 2019-2133. \nolinkurl{10.1093/rfs/hhy131}. \emph{Tier 1. Read: abstract.} \textbf{For us:} results depend on how much weight illiquid names get. Liquid-only tests may be the honest ones.
\item Jensen, T., B. Kelly, and L. Pedersen. "Is There a Replication Crisis in Finance?" Journal of Finance 78(5), 2465-2518. \nolinkurl{10.1111/jofi.13249}. \emph{Tier 1. Read: abstract.} \textbf{For us:} treat our 53 tests as one or a few themes. A shrinkage model suits correlated tests better than counting passes.
\item Harvey, C., and Y. Liu. "False (and Missed) Discoveries in Financial Economics." Journal of Finance 75(5), 2503-2553. \nolinkurl{10.1111/jofi.12951}. \emph{Tier 1. Read: abstract.} \textbf{For us:} a failed test is not proof of no effect. We should report our chance of a miss, not only of a false alarm.
\item Goyal, A., I. Welch, and A. Zafirov. "A Comprehensive 2022 Look at the Empirical Performance of Equity Premium Prediction." Review of Financial Studies 37(11), 3490-3557. \nolinkurl{10.1093/rfs/hhae044}. \emph{Tier 1. Read: RePEc abstract.} \textbf{For us:} this is the closest match to our exercise, a time-series predictor over a few market-wide episodes. Our result is common in this setting.
\item Inoue, A., and L. Kilian. "In-Sample or Out-of-Sample Tests of Predictability: Which One Should We Use?" Econometric Reviews 23(4), 371-402. \nolinkurl{10.1081/etc-200040785}. \emph{Tier 1. Read: RePEc abstract.} \textbf{For us:} a failed out-of-sample split can reflect low power alone. The 2019 split should not be our only test.
\item Hjalmarsson, E. "Predicting Global Stock Returns." Journal of Financial and Quantitative Analysis 45(1), 49-80. \nolinkurl{10.1017/s0022109009990469}. \emph{Tier 1. Read: abstract of the IFDP version.} \textbf{For us:} pooling many contracts is a known route to more power, but standard panel inference is unsafe here.
\item Boudoukh, J., M. Richardson, and R. Whitelaw. "The Myth of Long-Horizon Predictability." Review of Financial Studies 21(4), 1577-1605. \nolinkurl{10.1093/rfs/hhl042}. \emph{Tier 1. Read: abstract of NBER working paper 11841.} \textbf{For us:} a position-based signal is persistent, so 4, 13 and 26-week results are not separate evidence. Rising strength with horizon proves little.
\end{enumerate}

\end{document}
